\documentclass[10pt,a4paper]{amsart}
\usepackage[margin=19mm]{geometry}
\usepackage{amsmath,amssymb,mathtools}
\usepackage[scr=rsfs,cal=euler]{mathalfa}
\usepackage{xfrac}
\usepackage[unicode,hidelinks,bookmarks=false]{hyperref}
\newcommand{\ie}{i.e.\ }
\newcommand{\cf}{cf.\ }
\newcommand{\resp}{respectively\ }
\newcommand{\Subsection}{\subsection}
\providecommand{\nobreakdash}{\nobreak\hbox{-}}
\setlength{\emergencystretch}{2em}
\multlinegap0pt
\usepackage{amssymb}
\usepackage{tikz}
\newtheorem{thm}{Theorem}
\title{Associated primes, witnesses, and omega invariants of monomial ideals}
\author{Jacob Miller, Mehrdad Nasernejad, \c{S}tefan O. Toh\v{a}neanu}
\date{Source: arXiv:2609.11407v1 (CC BY 4.0)}
\begin{document}
\maketitle

\noindent\emph{Source and licence.} Associated primes, witnesses, and omega invariants of monomial ideals. Version arXiv:2609.11407v1 (CC BY 4.0), distributed by arXiv under the Creative Commons Attribution 4.0 International licence, http://creativecommons.org/licenses/by/4.0/, as recorded in the arXiv licence metadata for this exact version and shown on the abstract page https://arxiv.org/abs/2609.11407v1. Full licence notice: \texttt{licenses/arxiv-2609.11407-CC-BY-4.0.txt}.\par\smallskip

\noindent\textit{Editorial scope.} Counted unit: the article's thm thm:matrix\_power\_witness (source0.tex lines 1337--1346). The article's complete original proof of it is delivered with it (source0.tex lines 1348--1420, one \texttt{proof} environment of 73 lines). No mathematics is written, repaired or re-derived: the statement and the proof are the article's own lines, in the article's own notation, and no retained line is substituted or re-flowed.  No further source-local context is needed: every cross-reference of every retained line resolves inside the two delivered spans.  Nothing was omitted from any retained span, so the package carries no omission record.  No retained line contains a citation, so the wrapper carries no bibliography and no bibliography entry.  No retained line contains a figure environment, an image-inclusion command or drawing code, so no image is delivered and the package declares no asset dependency.  Editorial formatting only, in addition: an AMS \texttt{amsart} preamble that restates the article's own declarations and macros used by the retained lines, namely the environments \texttt{thm} on separate counters, together with the standard packages those lines need.  The retained environments print in this document as Theorem 1 in order of appearance, whereas the article prints the same items with the numbering of its own full document.  The complete native source of the article as distributed in the arXiv e-print for this exact version is bundled unchanged as \texttt{sources/210018/source0.tex}.
\begin{thm}{\rm (Matrix witness criterion for powers)}  \label{thm:matrix_power_witness}
%\leavevmode\newline
Let $I \subset R=\mathbb K[x_1,\ldots,x_k]$ be a monomial ideal minimally generated by $G(I)=\{f_1,\ldots,f_s\}$, and with exponent matrix $A$. Fix $n\geq1$ and let
$S=\{i_1,\ldots,i_m\}\subseteq[k].$  Suppose there exist vectors
 $\lambda^{(1)},\ldots,\lambda^{(m)} \in\mathbb N_0^s$ with $|\lambda^{(u)}|=n$ for all $1\leq u \leq m$
 such that the associated monomials form a compatible $S$-dominant witness system. Set
 $t_u=(A\lambda^{(u)})_{i_u}.$ Assume, in addition, that $t_u\geq 1$ for all $1\leq u \leq m$.
 If the integer system $\mathcal E_{S,n}$ has no solution, then
 $P_S=\langle x_{i_1},\ldots,x_{i_m}\rangle \in \operatorname{Ass}(R/I^n).$
\end{thm}

\begin{proof}
 For each $u=1,\ldots,m$, let $g_u=f^{\lambda^{(u)}}\in I^n.$ The exponent vector of $g_u$ is $b^{(u)}$.
For each $j\notin S$, set
\[c_j:=\max\left\{(A\lambda)_j:\lambda\in\mathbb N_0^s,\ |\lambda|=n\right\}.\]
This maximum exists because there are only finitely many vectors $\lambda\in\mathbb N_0^s$ satisfying $|\lambda|=n$.
 Define the monomial
\begin{equation}
M
=
\prod_{u=1}^m x_{i_u}^{t_u-1}
\prod_{j\notin S}x_j^{c_j}.
\label{eq:matrix_witness_M}
\end{equation}
The assumption $t_u\geq1$ guarantees that all exponents in \eqref{eq:matrix_witness_M} are nonnegative.

We first prove that $P_S\subseteq(I^n:M).$ Fix $u\in\{1,\ldots,m\}$. 

We claim that $g_u\mid x_{i_u}M.$ With the previous notation that for an arbitrary monomial $h$, $d_i(h)$ denotes the exponent of the variable $x_i$ showing up in $h$, we have $d_{i_u}(g_u)=t_u=d_{i_u}(x_{i_u}M)$.
For $v\neq u$, compatibility gives that $
d_{i_v}(g_u)
=
b^{(u)}_{i_v}
<
b^{(v)}_{i_v}
=
t_v.
$
Since all exponents are integers, we can deduce that $
b^{(u)}_{i_v}
\leq
t_v-1
=
d_{i_v}(M).
$
Finally, for $j\notin S$, we have $
d_j(g_u)
\leq
c_j
=
d_j(M).
$
Hence, $g_u\mid x_{i_u}M.$ Therefore, $x_{i_u}M\in I^n,$
and consequently $x_{i_u}\in(I^n:M).$ But this holds for every $u$, so we obtain
$P_S\subseteq(I^n:M).$

We now prove the reverse inclusion. Let $N$  be an arbitrary monomial such that  $N\notin P_S$.
 Then none of  $x_{i_1},\ldots,x_{i_m}$ divides $N$. Hence, we have $d_{i_u}(N)=0$
 for every $u$, and so we can deduce that
\begin{equation}
d_{i_u}(NM)
=
t_u-1
\qquad
\text{for every }u.
\label{eq:degree_NM}
\end{equation}
Suppose, to the contrary, that $NM\in I^n.$ Since $I^n$ is generated by products of $n$ generators of $I$, there
exists $\lambda\in\mathbb N_0^s$ with $|\lambda|=n$ such that  $f^\lambda\mid NM.$
It follows from \eqref{eq:degree_NM} that
\[
(A\lambda)_{i_u}
\leq
t_u-1
<
t_u
\qquad
\text{for every }u.
\]
Thus,  $\lambda$ is a solution of the escape system $\mathcal E_{S,n}$, contradicting the assumed infeasibility of that
system. Therefore, $NM\notin I^n$ for every monomial $N\notin P_S$. Hence, $(I^n:M)\subseteq P_S.$
Together with $P_S\subseteq(I^n:M)$, this leads to  $(I^n:M)=P_S.$
 We therefore get   $P_S\in\operatorname{Ass}(R/I^n)$, as required.
\end{proof}

\end{document}
